% ─── Research Questions and Scope ────────────────────────────────────────────
% Target: 1 to 2 pages.

\label{sec:rq}

% ─── Merged Research Questions ───────────────────────────────────────────────
% Single reconciled RQ set. The introduction (introduction.tex) and this section
% (research_questions.tex) pose the SAME numbered RQ1/RQ2/RQ3; the literature
% review (literature_review.tex, sec:lit:gap) states the same content as THREE
% GAPS, not a second numbered set, and maps gap-for-gap onto these questions
% ("...it is this that motivates Research Question 2"). The three gaps and the
% three RQs are therefore one set, restated here with each gap folded into the
% question that closes it.

This study is organised around three research questions, each posed as the answer to one
of the three gaps identified in Section~\ref{sec:lit:gap}: the missing NUTS3-resolved
demand basis for a continental assessment in the updated REPowerEU context (gap one),
the absence of a hydrogen-versus-heat-pump comparison on a consistent operating-cost
basis that bounds hydrogen's share by where it can win and caps each scenario at or
below that bound (gap two),
and the severed link between the heating choice and the power sector that electrified
heat runs on, of which the import-reliant Swiss case is the sharpest instance (gap
three).

RQ1 closes gaps one and two together at the EU scale; RQ2 takes up gap three for
Switzerland as a full national case in its own right; and RQ3 names the
buildings-to-power coupling that the severed link implies. Each question is carried in
full so that the reader can hold the analysis against the question it sets out to
settle. RQ1 and RQ2 are answered in full by the analysis that follows; RQ3 is scoped to
the buildings-to-power synergy, which this paper prices on the same operating-cost basis
as the rest of the analysis, while the cross-sector synergies with industry and
transport are treated qualitatively and deferred to future work, the precise boundary
being set out in the RQ3 statement and the scope subsection below.

\begin{description}
  \item[RQ1 (EU, UK and Switzerland, country-level).] Under what conditions is
  decarbonised hydrogen cost-competitive for building heat by 2050 under a net-zero
  2050 requirement, where do those conditions hold, and how do the resulting pathways
  compare on economic and environmental cost bases? RQ1 takes up the first two gaps together. It addresses hydrogen's heating role at NUTS3 resolution across the EU-27, the United Kingdom
  and Switzerland, and tests the hydrogen-versus-heat-pump trade-off on a
  consistent operating-cost basis, the trade-off the critical literature finds
  decisive but parameter-sensitive \citep{Rosenow2022PipeDream,ICCT2023h2heating}. It asks where hydrogen is deployed, on what cost basis, and in competition with which alternative, and it answers by
  deriving hydrogen's heating share from where the fuel can win the cold-snap peak
  in a merit-order dispatch instead of imposing it as an exogenous target, reading
  levelised cost of heat, carbon abatement and infrastructure cost off the same
  modelled pathway.

  \item[RQ2 (Switzerland, national case).] How can Switzerland make use of the EU's
  adoption of hydrogen-based technologies to contribute to its own net-zero 2050
  aim? RQ2 closes the third gap for Switzerland, a high-cost, hydro-rich, cavern-poor system
reliant on imported hydrogen \citep{Gabrielli2020,IEA2023switzerland}, carried through the same framework as a
  full national case, neither subordinate to the EU-wide assessment nor split off
  into a separate study, because its distinctive position, a near-decarbonised
  grid, no domestic salt-cavern storage \citep{Caglayan2020caverns}, a reliance on
  imported hydrogen whose delivered price is the principal lever, and a
  winter-electricity adequacy gap that any firm-capacity role would have to address
  \citep{Mellot2024Mitigating}, warrants investigation on its own terms and on the
  same operating-cost basis as the rest of the analysis. We examine the sensitivity
  of the optimal Swiss buildings pathway to the import price across the central,
  rapid and stranded supply trajectories, and ask under what conditions, if any,
  hydrogen earns a place in Swiss residential heat.

  \item[RQ3 (cross-sector synergy, scoped to the power sector).] What synergies can
  be established between hydrogen in the buildings sector and other sectors? RQ3
  names the buildings-to-power coupling that the severed link of the third gap
  implies. Buildings do not decarbonise in isolation. The electrified heat load
  reshapes power demand, and any hydrogen used for heat draws on the same
  production, storage and transport system that serves industry and transport. This
  paper answers in full the one synergy that follows most directly from its own
  results, the coupling between buildings-sector hydrogen and the power sector,
  carrying the electrified heat load into the electricity merit order and assessing
  hydrogen's role in the firm-capacity stack on the same operating-cost basis
  (Section~\ref{sec:power}), the role the system-value literature locates in scarce
  peak hours, not in annual energy \citep{Brown2018Synergies,Zeyen2021Mitigating}. The wider synergies with industry and
  transport are treated qualitatively in the Discussion
  (Section~\ref{sec:discussion}); a full cross-sector optimisation lies beyond the
  present scope and is left to subsequent work.
\end{description}

\subsection{Research Scope and Methodological Boundaries}

The following scope statement defines the geographic, temporal, and technological
boundaries within which RQ1, RQ2, and RQ3 are addressed. The analysis covers
residential space heating and hot water demand across 29 countries (the EU-27,
the United Kingdom, and Switzerland) at NUTS3 regional resolution. The modelling
period spans 2025 to 2050 at milestone years (2025, 2030, 2040, 2050). Three
building types are distinguished, single-family homes, multi-family dwellings
(including apartment blocks), and other residential buildings, reflecting
differences in technology feasibility and retrofit economics. Eight heating
technologies compete in each region: gas boiler, oil boiler, biomass boiler,
resistance heater, air-source heat pump, ground-source heat pump, district
heating, and hydrogen boiler. To this building-level competition the study adds
two further layers that test hydrogen beyond the boiler comparison: the
district-heating supply stack and the power sector on which electrified heat
depends. The power layer commits the firm fleet with an explicit unit-commitment
optimisation (ramp, minimum stable output, reserve and start-up constraints).

Three boundaries delimit the study. First, it treats the residential sub-sector
only. Commercial and public buildings, industrial process heat, and end uses
other than space heating and hot water are excluded. Second, it stops at the
meter and the wholesale market for the buildings question itself, with network
reinforcement entering through an explicit infrastructure-cost account, short of a nodal grid model. Third, and following directly from RQ3 above, sector
coupling beyond the buildings-to-power link is not modelled endogenously.
Synergies between hydrogen in buildings and its use in industry or transport are
discussed but not co-optimised, and are flagged throughout as a direction for
future work. This last boundary is a deliberate design choice, not an omission.
The buildings-to-power link is the synergy that follows directly from the
paper's own dispatch results and can be priced on the same operating-cost basis. A full industry-and-transport co-optimisation would require a separate
model and is held over and not treated superficially here. Within these bounds the paper aims to price hydrogen for heat the
way a system operator actually would, and to let its role be decided by that
test, not assumed.
